\section{Introduction}

Telomeres are the nucleoprotein structures that cap the ends of linear eukaryotic chromosomes, where they solve the end-replication problem and prevent the chromosome terminus from being recognised as a double-strand break \cite{delange2005shelterin,palm2008howshelterin}. In somatic human cells, progressive loss of terminal TTAGGG repeats with each division eventually triggers replicative senescence, linking telomere attrition to both cellular and organismal aging \cite{harley1990telomeres,blackburn2006human}. The rate and pattern of telomere shortening are exquisitely sensitive to the balance between proliferation and telomere maintenance, and a small but clinically important class of inherited telomere biology disorders (TBDs) demonstrates that deviations from this balance produce diseases of the bone marrow, lung and skin \cite{armanios2012telomeres,bertuch2016molecular}. As a result, telomere length is widely used as a molecular biomarker of replicative history, genomic instability and disease risk \cite{blackburn2006human,aubert2012telomere}.

Despite decades of interest, the methods used to measure telomere length remain a bottleneck. Terminal restriction fragment (TRF) Southern blots return a bulk mean from a smear on a gel \cite{aubert2012telomere}, flow-FISH and qPCR measure relative telomeric signal rather than individual telomere lengths \cite{baerlocher2006flow,cawthon2002telomere}, and PCR-based assays such as STELA and TeSLA that do attempt to resolve individual telomeres are limited by polymerase processivity and are preferentially biased toward the shortest fraction of the distribution \cite{baird2003extensive,lai2018comparison}. Short-read sequencing approaches such as TelSeq extract a relative telomere content from whole-genome data but cannot recover the length of any single telomere \cite{ding2014telseq}. Long-read sequencing technologies together with the recently completed telomere-to-telomere (T2T) human reference genome now make it feasible, in principle, to align and measure individual intact telomeric reads \cite{nurk2022complete,tham2023comprehensive,schmidt2023targeted,karimian2024human}, but no pipeline has yet combined these advances into a clinical workflow capable of resolving the full length distribution and using it to distinguish healthy aging from disease.

We close this gap with a digital telomere measurement pipeline built around nanopore long-read sequencing. A biotinylated oligonucleotide complementary to the 3$'$ overhang, combined with restriction digestion of non-telomeric genomic DNA, enriches telomeric reads by several thousand-fold without distorting the length distribution. A bioinformatic tool, Telometer, aligns telomeric repeats to chromosomal termini of the T2T reference and defines the telomere boundary as the span between the first terminal repeat and the last two consecutive TTAGGG units before the subtelomere. The resulting per-read lengths provide a chromosome-anchored view of every intact telomere in the sample. We validate the method against TRF ($n{=}14$) and flow-FISH ($n{=}7$), quantify its precision by bootstrapping, and apply it to genetically defined hESC models, peripheral blood from 14 healthy donors spanning ages 18 to 77, peripheral blood and bone marrow from patients with RTEL1 or TINF2 variants, and matched benign and malignant colon tissue from 20 colorectal carcinoma patients \cite{walne2013constitutional,tummala2015poly,savage2008tinf2,alder2018diagnostic}.

Our main contributions are as follows. (i) We establish that digital telomere measurement plateaus at a precision of 30--40 bp, correlates tightly with TRF and flow-FISH, and that TRF systematically overestimates mean telomere length by 1000--3000 bp. (ii) We show that the pipeline recovers $\sim$40 bp/day of telomere attrition in TERT-null hESCs and $\sim$1000 bp of de novo elongation after three days of telomerase overexpression, with chromosome-anchored evidence that telomerase preferentially extends the chromosomes with shortest baseline telomeres \cite{hemann2001shortest,britt2012preferential,teixeira2004intertelomere}. (iii) In peripheral blood leukocytes from 14 healthy donors, mean and median telomere length decline by $\sim$24 bp/year, and the third quartile shortens faster than the first---an asymmetry that is abolished in RTEL1 carriers and in TERT-null hESCs. (iv) A simple machine-learning classifier trained on the length distribution and donor age distinguishes healthy donors from symptomatic TBD patients (AUC$=$0.98), from asymptomatic carriers that flow-FISH cannot detect (AUC$=$0.93), and from the combined group (AUC$=$0.91). (v) In matched benign and malignant colorectal tissue, tumour telomeres are shorter than benign epithelium in 75\% of the 20 patients and show no correlation with donor age, consistent with telomerase-stabilised tumour distributions \cite{barthel2017systematic,roka2023tumor}.

\section{Related Work}

\paragraph{Analog telomere measurement.} The canonical methods for measuring telomere length are TRF Southern blotting, flow-FISH, Q-FISH and qPCR. Each delivers a coarse mean or relative content rather than individual telomere lengths \cite{aubert2012telomere,baerlocher2006flow,cawthon2002telomere,lai2017comparison}. TRF additionally suffers from incomplete digestion of subtelomeric DNA, which inflates the apparent mean, and qPCR reports only a ratio normalised to a single-copy gene \cite{cawthon2002telomere}. Flow-FISH, the current clinical standard for PBLs, is limited to mean telomere content in PBLs and struggles to distinguish asymptomatic TBD carriers from healthy individuals \cite{alder2018diagnostic}.

\paragraph{PCR-based single-telomere assays.} STELA was the first method to resolve individual chromosome ends by ligating a linker to the G-overhang and amplifying a tagged terminal fragment, revealing extensive allelic variation and ultrashort telomeres at senescence \cite{baird2003extensive}. TeSLA extended this idea to the genome as a whole by restriction digesting, ligating adapters and amplifying short telomeric fragments, making the assay sensitive to the shortest fraction of the distribution \cite{lai2018comparison}. Both methods are valuable for the shortest-telomere fraction but are limited by polymerase processivity and by systematic under-representation of long telomeres.

\paragraph{Sequencing-based estimates of telomere content.} Short-read sequencing approaches such as TelSeq estimate relative telomere content by counting reads carrying telomeric repeats in whole-genome data, and have been applied to large cohorts such as TCGA to relate telomere content to cancer biology \cite{ding2014telseq,barthel2017systematic,demanelis2020determinants}. These approaches do not recover individual telomere lengths and cannot separate the interquartile structure of the length distribution from a change in total telomeric content.

\paragraph{Long-read sequencing of telomeres.} PacBio HiFi and Oxford Nanopore sequencing have begun to offer single-molecule, per-chromosome telomere measurements \cite{tham2023comprehensive,schmidt2023targeted,karimian2024human,heller2019telomeres}. Alignment to the T2T reference provides, for the first time, unambiguous anchoring of individual telomeres to specific chromosome arms \cite{nurk2022complete}. However, published pipelines have either been limited to whole-genome sequencing (which is costly per sample) or to highly targeted capture platforms, and none has combined chromosome-anchored measurements of the complete length distribution with clinical applications in aging and disease diagnosis.

\paragraph{Telomere biology disorders and the clinical pathway.} Dyskeratosis congenita and related TBDs are driven by germline mutations in components of the telomerase holoenzyme (TERT, TERC, PARN, DKC1) or the shelterin complex (TINF2) and in the helicase RTEL1 \cite{tummala2015poly,walne2013constitutional,savage2008tinf2,armanios2012telomeres,bertuch2016molecular}. Clinical triage combines flow-FISH, population norms and targeted exon sequencing \cite{alder2018diagnostic}. Because TBDs show genetic anticipation and because asymptomatic carriers may have telomere lengths within the lower quintile of the population, a higher-resolution assay capable of detecting carrier status would fill a diagnostic gap---a gap we address here.

\paragraph{Telomere homeostasis and length regulation.} Telomerase preferentially elongates the shortest telomeres in yeast and in mammalian systems \cite{teixeira2004intertelomere,hemann2001shortest,britt2012preferential}, and shelterin components regulate telomerase access and act as a negative-feedback rheostat on length \cite{palm2008howshelterin,hockemeyer2015control}. In cancer, telomerase reactivation (often via TERT promoter mutations) stabilises tumour telomere length distributions at short but persistent set-points \cite{chiba2015cancer,barthel2017systematic}. Our single-molecule data allow these rheostat-level regulatory features to be read out from the interquartile structure of the distribution itself.
